\section{Introducción}

La innovación empresarial requiere recursos para investigación y desarrollo, adquisición de maquinaria, software, ingeniería, capacitación y otras actividades. El volumen de esa inversión importa, pero también importan su fuente de financiación y su destino. Una empresa que depende de recursos propios enfrenta restricciones distintas a otra que accede a crédito, capital privado o instrumentos públicos. De la misma forma, la compra de maquinaria cumple una función diferente a la inversión en I+D o en activos digitales.

Este informe caracteriza tres dimensiones de la inversión empresarial en actividades científicas, tecnológicas y de innovación (ACTI): su valor como porcentaje de los ingresos o ventas, la distribución del valor financiado entre fuentes y la distribución de la inversión entre usos. Se utilizan la Encuesta de Desarrollo e Innovación Tecnológica en la industria manufacturera (EDIT X, 2019--2020) y la Encuesta de Desarrollo e Innovación Tecnológica en los sectores de servicios y comercio (EDITS VIII, 2020--2021).

El análisis toma 2019 para manufactura y 2021 para servicios y comercio. En manufactura se selecciona 2019 para evitar que el choque excepcional de la pandemia determine la caracterización de la inversión. En servicios y comercio se utiliza 2021, el último año disponible. Los valores monetarios están expresados en miles de pesos corrientes en los microdatos. Las participaciones y la inversión sobre ventas no dependen de esa escala monetaria, pero la diferencia de años y universos debe mantenerse al interpretar los resultados.

\section{Fuente, medidas y alcance}

La base de manufactura contiene 6.797 observaciones empresariales y la base de servicios y comercio contiene 8.812. Para calcular la intensidad agregada se sumó la inversión en ACTI de las empresas con ventas positivas y se dividió por la suma de sus ventas nacionales y exportaciones. Este cociente entre sumas mide qué proporción de las ventas del conjunto se destinó a ACTI; no es el promedio simple del porcentaje de cada empresa.

Las fuentes de financiación se clasifican en recursos propios, empresas del grupo, recursos públicos, banca privada, otras empresas, capital privado y cooperación o donaciones. Los usos incluyen I+D interna y externa, maquinaria y equipo, TIC, software y datos, mercadotecnia, propiedad intelectual, asistencia técnica y consultoría, ingeniería y diseño, formación y capacitación, edificaciones y métodos organizativos.

Los controles contables muestran que, en las dos encuestas, la suma de las fuentes coincide con el total de financiación reportado y la suma de los usos coincide con la inversión total en ACTI. Esta identidad se cumple tanto en el valor agregado como en cada observación. Por lo tanto, las participaciones presentadas en las figuras suman 100\% y no requieren un ajuste residual.

\begin{table}[H]
\centering
\caption{Volumen e intensidad de la inversión empresarial en ACTI}
\label{tab:resumen}
\begin{threeparttable}
\begin{tabular}{lrr}
\toprule
 & Manufactura, 2019 & Servicios y comercio, 2021 \\
\midrule
Observaciones empresariales & 6.797 & 8.812 \\
Empresas con inversión positiva & 1.307 & 1.865 \\
Participación con inversión positiva & 19,2\% & 21,2\% \\
Ventas (billones de pesos corrientes) & 270,3 & 684,9 \\
Inversión ACTI (billones de pesos corrientes) & 2,20 & 4,40 \\
Inversión ACTI / ventas & 0,81\% & 0,64\% \\
Mediana ACTI / ventas entre inversoras & 0,88\% & 0,47\% \\
\bottomrule
\end{tabular}
\begin{tablenotes}\footnotesize
\item Nota: el cociente agregado restringe el denominador a empresas con ventas positivas. Los valores monetarios se convierten de miles de pesos a billones de pesos. La mediana se calcula entre empresas con inversión positiva y ventas positivas.
\item Fuente: cálculos del CJC con microdatos DANE, EDIT X y EDITS VIII.
\end{tablenotes}
\end{threeparttable}
\end{table}

\section{La inversión representó menos del 1\% de las ventas}

La Figura \ref{fig:intensidad} muestra que la inversión en ACTI representó 0,81\% de las ventas en la industria manufacturera y 0,64\% en servicios y comercio. La diferencia fue de 0,17 puntos porcentuales. Dado que las encuestas cubren sectores y años distintos, la magnitud de cada cociente es más informativa que su orden: en ambos casos la inversión fue inferior a un peso por cada cien pesos de ventas.

\begin{figure}[H]
\centering
\includegraphics[width=0.96\textwidth]{informe_edit_intensidad.png}
\caption{Inversión empresarial en ACTI como porcentaje de las ventas}
\label{fig:intensidad}
\end{figure}

La inversión positiva se concentró en cerca de una quinta parte de las observaciones: 19,2\% en manufactura y 21,2\% en servicios y comercio. Entre esas empresas, la mediana de la inversión sobre ventas fue 0,88\% en manufactura y 0,47\% en servicios y comercio. La diferencia entre el cociente agregado y la mediana empresarial indica que el resultado depende de la distribución del tamaño y de la intensidad de inversión entre empresas; el agregado no describe a una empresa representativa.

\section{Los recursos propios fueron la principal fuente}

La Figura \ref{fig:fuentes} reporta la composición del valor financiado. Los recursos propios representaron 80,1\% en manufactura y 76,8\% en servicios y comercio. La banca privada fue la segunda fuente en manufactura, con 14,3\%, mientras los recursos públicos ocuparon ese lugar en servicios y comercio, con 12,3\%.

\begin{figure}[H]
\centering
\includegraphics[width=0.98\textwidth]{informe_edit_fuentes.png}
\caption{Fuentes de financiación de la inversión empresarial en ACTI}
\label{fig:fuentes}
\end{figure}

Los recursos públicos representaron 0,4\% de la financiación manufacturera. La diferencia frente al 12,3\% de servicios y comercio es amplia, pero la encuesta no identifica su mecanismo. Puede reflejar diferencias en los instrumentos disponibles, en la propensión de las empresas a solicitarlos, en los criterios de selección o en la composición de actividades incluidas en cada operación. Distinguir entre esas explicaciones requiere información sobre solicitudes, aprobación y características de los proyectos.

La suma de recursos propios, empresas del grupo, banca privada, otras empresas y capital privado representó 99,6\% de la financiación manufacturera y 86,9\% de la financiación de servicios y comercio. La cooperación y las donaciones tuvieron una participación inferior a 1\% en ambos casos. La financiación de ACTI, por lo tanto, dependió principalmente de fuentes privadas y, dentro de ellas, de la liquidez de las propias empresas.

\section{Los usos reflejan diferencias en la actividad productiva}

La Figura \ref{fig:usos} muestra que la manufactura destinó 49,3\% de la inversión a maquinaria y equipo. La I+D interna representó 15,1\% y las TIC, software y datos 7,7\%. Estas tres categorías reunieron 72,1\% de la inversión manufacturera.

\begin{figure}[H]
\centering
\includegraphics[width=0.98\textwidth]{informe_edit_usos.png}
\caption{Usos de la inversión empresarial en ACTI}
\label{fig:usos}
\end{figure}

En servicios y comercio, los tres usos principales tuvieron pesos más cercanos: maquinaria y equipo representó 26,1\%, TIC, software y datos 25,4\%, e I+D interna 22,0\%. En conjunto sumaron 73,6\%. La principal diferencia frente a la manufactura fue la menor concentración en maquinaria y el mayor peso de activos digitales. La I+D externa también tuvo una participación mayor en servicios y comercio: 8,1\%, frente a 2,7\% en manufactura.

La composición de los usos es consistente con diferencias tecnológicas entre actividades, pero no permite concluir que un destino produzca más innovación o productividad que otro. La EDIT y la EDITS reportan cuánto se invierte en cada actividad; evaluar su retorno requiere relacionar estos usos con resultados posteriores de ventas, productividad, nuevos productos o procesos.

\section{Conclusiones}

La inversión empresarial en ACTI representó menos del 1\% de las ventas en los dos universos estudiados. La manufactura registró una intensidad agregada de 0,81\% en 2019 y los servicios y el comercio de 0,64\% en 2021. En ambos casos, cerca de una quinta parte de las observaciones reportó inversión positiva.

Los recursos propios financiaron más de tres cuartas partes de la inversión. Esta dependencia fue mayor en manufactura, donde alcanzó 80,1\%. Los recursos públicos tuvieron un peso mucho mayor en servicios y comercio que en manufactura. El resultado describe la composición de la financiación, pero no establece si el apoyo público genera inversión adicional ni por qué llega con distinta intensidad a los dos universos.

Los usos también fueron distintos. La manufactura concentró 49,3\% de la inversión en maquinaria y equipo. En servicios y comercio, la inversión se distribuyó de forma más pareja entre maquinaria, TIC e I+D interna. La pregunta que queda abierta es si estas diferencias de financiación y destino se traducen en resultados distintos de innovación y productividad, y qué empresas quedan por fuera de la inversión por depender principalmente de sus propios recursos.

\section*{Fuentes y reproducibilidad}
\addcontentsline{toc}{section}{Fuentes y reproducibilidad}

Los cálculos usan los archivos \texttt{Datos/Raw/EDIT/EDIT\_X\_2019\_2020.dta} y \texttt{Datos/Raw/EDIT-S/Estructura\_EDITS\_VIII\_2020\_2021.dta} del repositorio CJC-Monitor. El script reproducible es \texttt{Code/03.2\_Informe\_EDIT\_FuentesUsos.py}. Las tablas agregadas que respaldan el informe se conservan en \texttt{Paper/tables}.

La página del DANE sobre la EDIT identifica estas dos operaciones como las últimas ediciones publicadas para industria manufacturera y para servicios y comercio: \url{https://www.dane.gov.co/index.php/estadisticas-por-tema-2/tecnologia-e-innovacion/encuesta-de-desarrollo-e-innovacion-tecnologica-edit}.
